\subsection{Regular subcomplexes and a derived obstruction}

The termwise lattices in \eqref{wo:eq:nonregular-face} do not form a
complex.  There are two natural ways to retain them: restrict the source
until its differential is regular, or enlarge the target by the missing
images.  These constructions give the same derived object.  For the
Wick--Ore differential, that object has torsion which is absent from the
pure ordered regular bar.

Continue with \(R=k[[z]]\), \(K=k((z))\), and the algebraic state
modules of Section~\ref{wo:sec}.  In particular, every vector uses
finitely many state monomials.  Only its coefficients are power series.
For a cochain complex \((C,d)\) of \(K\)-vector spaces, let
\(L^n\subset C^n\) be prescribed \(R\)-submodules spanning
\(C^n\) over \(K\).  They need not be preserved by \(d\).  Define
\begin{equation}\label{woreg:eq:closures}
 C_-^n=L^n\cap d^{-1}(L^{n+1}),\qquad
 C_+^n=L^n+dL^{n-1}.
\end{equation}
Both differentials are restrictions of \(d\).

\begin{lemma}[Two regular subcomplexes]\label{woreg:lem:closures}
The complex \(C_-\) is the largest subcomplex contained in \(L\),
and \(C_+\) is the smallest subcomplex containing \(L\).
The inclusion \(C_-\hookrightarrow C_+\) is a quasi-isomorphism.
These constructions are functorial for chain maps preserving the
prescribed graded submodules.
\end{lemma}
\begin{proof}
If \(x\in C_-^n\), then \(dx\in L^{n+1}\) and
\(d(dx)=0\in L^{n+2}\), so \(dx\in C_-^{n+1}\).
The identity \(d^2=0\) also makes \(C_+\) a subcomplex.
Their extremal properties follow directly from
\eqref{woreg:eq:closures}.

Every element of \(C_+^n\) has the form \(l+dm\), with
\(l\in L^n\) and \(m\in L^{n-1}\).  If its class is a cycle
in \(C_+/C_-\), then \(dl\in C_-^{n+1}\subset L^{n+1}\).
Thus \(l\in C_-^n\), and the quotient class of \(l+dm\) is the
boundary of the class of \(m\).  The quotient is acyclic, proving
the quasi-isomorphism.  A chain map preserving \(L\) preserves both
the inverse-image condition and the sum in
\eqref{woreg:eq:closures}.
\end{proof}

For the reduced two-point complex put
\[
 N_R=(R\otimes_k W\otimes_k1)
       \oplus(R\otimes_k1\otimes_k W),\qquad
 N_K=K\otimes_RN_R.
\]
Its prescribed lattices and differential are
\[
 X_R=M_R\oplus M_R,\qquad Y_R=N_R\oplus M_R,\qquad
 d(p,q)=E(\Theta_{12}p+\Theta_{21}q).
\]
Write \(X_K=K\otimes_RX_R\) and \(Y_K=K\otimes_RY_R\).
Here \(X_R\) has degree \(-2\), \(Y_R\) has degree \(-1\),
and the displayed differential takes values in the mixed summand
of \(Y_K\).  Formula~\eqref{woreg:eq:closures} gives
\begin{align}
 C_-^{-2}&=X_R\cap d^{-1}(M_R),&
 C_-^{-1}&=N_R\oplus M_R,\label{woreg:eq:lower}\\
 C_+^{-2}&=X_R,&
 C_+^{-1}&=N_R\oplus(M_R+EM_R).\label{woreg:eq:upper}
\end{align}
The ordered operators preserve \(M_R\), and
\(d(X_R)=EM_R\), which explains the upper target.
Write
\begin{equation}\label{woreg:eq:defect}
 D_E=M_R/(M_R\cap EM_R).
\end{equation}

\begin{theorem}[Torsion in the regular complexes]\label{woreg:thm:torsion}
The inclusion \(C_-\hookrightarrow C_+\) induces the identifications
\[
 H^{-2}(C_-)\cong H^{-2}(C_+)\cong\gamma(M_R),\qquad
 H^{-1}(C_-)\cong H^{-1}(C_+)\cong N_R\oplus D_E.
\]
All other reduced cohomology groups vanish.  The module \(D_E\) is
nonzero and every element is killed by a power of \(z\).  For every
\(r,s\ge0\), it contains a submodule isomorphic to
\(R/z^{r+s+2}R\).  Consequently no single power of \(z\) kills
\(D_E\).  Neither regular complex is isomorphic in \(D(R)\) to
the pure ordered regular complex obtained by setting \(E=1\).
\end{theorem}
\begin{proof}
Theorem~\ref{wo:thm:homology} and regularity of \(G\) identify the
kernel on \(X_R\) with \(\gamma(M_R)\).  Every such cycle belongs
to \(C_-\), so the two degree-\(-2\) kernels agree.  The upper
mixed cokernel is
\[
 (M_R+EM_R)/EM_R\cong M_R/(M_R\cap EM_R).
\]
The image of the lower differential is exactly \(M_R\cap EM_R\).
Indeed, if \(x\) belongs to this intersection, then
\(a=\Theta_{12}^{-1}E^{-1}x\in M_R\) and
\(d(a,0)=x\).  This proves both cohomology formulas, including
their compatibility with the inclusion.

For \(x\in M_R\), the vector \(E^{-1}x\) has finite state
support and a finite coefficient pole order.  For some \(n\),
\(z^nE^{-1}x\in M_R\), hence
\(z^nx\in M_R\cap EM_R\).  Thus \(D_E\) is elementwise
\(z\)-power torsion.

Use the paired-jet plane of Proposition~\ref{wo:prop:lattice}, with
\(v=J_{1,r}^+\otimes J_{2,s}^-\),
\(t=\tau_1\otimes\tau_2\), \(p=r+s+2\), and
\(L_0=Rv\oplus Rt\).  Both \(E\) and \(E^{-1}\) preserve
\(V_0=Kv\oplus Kt\).  Since \(M_R\cap V_0=L_0\), it follows
that \(EM_R\cap V_0=EL_0\).  Therefore the natural map
\[
 L_0/(L_0\cap EL_0)\longrightarrow D_E
\]
is injective.  Formula~\eqref{wo:eq:lattice-sum} identifies its source
with \(R/z^pR\), generated by the class of \(v\).
Its annihilator is exactly \(z^pR\).  The values of \(p\) are
unbounded, proving the torsion assertions.

For \(E=1\), the differential already preserves the prescribed
lattices and surjects onto their mixed summand.  Its degree-\(-1\)
cohomology is \(N_R\), which is \(R\)-free.  The nonzero torsion
in \(N_R\oplus D_E\) rules out an isomorphism of these cohomology
modules and therefore an isomorphism in \(D(R)\).
\end{proof}

Both regular complexes become the original complex after tensoring
with \(K\).  For the lower source, a sufficiently high power of
\(z\) clears the finitely many poles of the differential of any
given finite state expression.  The upper target spans the same
\(K\)-module.  The derived obstruction is therefore a property of
these specified regular objects; it does not contradict the
meromorphic comparison.

On the paired-jet plane the two complexes have elementary matrices.
The operators \(O_{12}\) and \(O_{21}\) vanish there, so
\(\Theta_{12}\) and \(\Theta_{21}\) are identities.  Put
\(\alpha=(-1)^r(r+s+1)!/(r!s!)\) and \(t'=z^{-p}t\).
In the upper target basis \((v,t')\) and source basis
\((v_{12},t_{12},v_{21},t_{21})\), the differential is
\[
 d_+=\begin{pmatrix}1&0&1&0\\
                    \alpha&z^p&\alpha&z^p\end{pmatrix}.
\]
For the lower source take
\((v_{12}-v_{21},t_{12},t_{21},z^pv_{21})\), and use \((v,t)\)
in the target.  Then
\[
 d_-=\begin{pmatrix}0&0&0&z^p\\0&1&1&\alpha\end{pmatrix},
 \qquad
 J^{-2}=\begin{pmatrix}
  1&0&0&0\\0&1&0&0\\-1&0&0&z^p\\0&0&1&0
 \end{pmatrix},\qquad
 J^{-1}=\begin{pmatrix}1&0\\0&z^p\end{pmatrix}.
\]
Multiplication gives \(d_+J^{-2}=J^{-1}d_-\).
The quotient differential is multiplication by the unit \(\alpha\)
between two copies of \(R/z^pR\).  Row and column operations give
the invariant factors \(1,z^p\) for each differential.
This finite plane exhibits the torsion class; it is not asserted to
be stable under current translation.

\subsection{The common-coefficient allocation coalgebra}

The two regular complexes also retain a coalgebra structure when the
coefficient tensor is specified explicitly.  This uses a general
observation.  Suppose \(C\) is an algebraic coaugmented dg
\(K\)-coalgebra and its prescribed graded submodule \(L\) contains
\(R\cdot1\), has \(\varepsilon(L)\subseteq R\), and satisfies
\[
 \Delta(L)\subseteq L\otimes_RL
                 \subseteq C\otimes_KC.
\]
Comultiplication is a finite sum on each element.  Every
\(R\)-submodule of a \(K\)-vector space is torsion-free and thus
flat over the discrete valuation ring \(R\).  The displayed tensor
maps and their iterates are therefore injective.

\begin{lemma}\label{woreg:lem:coalgebra-closure}
Under these hypotheses \(C_+=L+dL\) is a coaugmented dg
\(R\)-coalgebra.  It is conilpotent if \(C\) is conilpotent.
\end{lemma}
\begin{proof}
The coderivation identity gives
\[
 \Delta(dx)=(d\otimes1+1\otimes d)\Delta(x)
       \in(dL\otimes_RL)+(L\otimes_RdL)
       \subseteq C_+\otimes_RC_+
 \qquad(x\in L),
\]
where the second summand uses the usual graded tensor sign.
Together with the hypothesis on \(\Delta(L)\), this proves closure.
The counit, coaugmentation and coassociativity identities restrict
from \(C\), since the corresponding tensor maps are injective.
The same injectivity makes any vanishing iterated reduced coproduct
remain zero in the subcoalgebra, proving conilpotence.
\end{proof}

For a general graded subcoalgebra lattice, this argument gives no
coalgebra structure on \(C_-\).  The two-point allocation model
has an additional support property which supplies that structure.

Define the two-point allocation coalgebra over the common field
\(K\) as follows.  Its coaugmentation is \(K\); its degree-\(-1\)
part is \(N_K\oplus M_K\), and its degree-\(-2\) part is
\(M_{12}\oplus M_{21}\).  Every degree-\(-1\) state is primitive.
The reduced coproduct of \(M_{12}\) is the corresponding tensor
of the first and second singleton states, and that of \(M_{21}\)
is the tensor in the opposite order.  These lie in independent
allocation summands.  Add the two coaugmentation terms to obtain
the full coproduct.  The differential is the row
\(E(\Theta_{12},\Theta_{21})\); it takes two-factor states to
primitive mixed one-factor states.  This is a dg coalgebra: its
reduced coderivation equation has zero on both sides, and the unit
cuts supply the full equation.  Every iterated reduced coproduct
of length three is zero.

\begin{proposition}\label{woreg:prop:regular-coalgebras}
Adjoin the unchanged coaugmentation \(R\) to \(C_-\) and
\(C_+\) in \eqref{woreg:eq:lower}--\eqref{woreg:eq:upper}.
Both are conilpotent dg \(R\)-coalgebras for these allocation cuts,
and their inclusion is a quasi-isomorphism of dg coalgebras.
The additional torsion classes in degree \(-1\) have primitive
mixed one-factor representatives.
\end{proposition}
\begin{proof}
The prescribed regular word modules form a graded subcoalgebra.
Lemma~\ref{woreg:lem:coalgebra-closure} therefore applies to \(C_+\).
For \(C_-\), every reduced cut of a degree-\(-2\) element lies in
the two regular singleton factors, both of which belong to \(C_-\).
All its other elements are primitive or coaugmented.  Thus \(C_-\)
is also closed under the coproduct.  The inclusion respects these
restrictions, and Lemma~\ref{woreg:lem:closures} proves its
quasi-isomorphism.  The representatives from
\eqref{woreg:eq:defect} lie in the primitive mixed one-factor part.
Their reduced coproduct vanishes already on chains; no K\"unneth
identification is required.
\end{proof}

The torsion in Theorem~\ref{woreg:thm:torsion} rules out even a zigzag
of dg-coalgebra quasi-isomorphisms from either regularization to the
pure ordered regular allocation coalgebra.  This assertion concerns
the explicitly defined common-coefficient tensor over \(R\).
Evaluation of a general configuration coalgebra at one configuration
has not been asserted to be monoidal.  A global regularization must
separately preserve configuration transitions, translations, simultaneous
collision divisors and the operations in higher arity.  Those conditions
remain part of the geometric recipient comparison.
